\section{Robustness to Frequently-Executed Tasks Sample Restriction}
\label{app:frequency_robustness}

The task sequences underlying our empirical analysis combine, for each occupation, the full set of O*NET tasks regardless of how often a task is actually performed on the job.
This raises the concern that rarely-executed tasks, which a worker might perform only rarely, lengthen the workflow and insert ``filler'' steps between the tasks that structure day-to-day production.
If they do, the average length of the AI chains we document in Prediction~\#1, the adjacency between AI-executed steps we exploit in Prediction~\#2, and the dispersion of AI-exposed steps we measure in Prediction~\#3 may partly reflect infrequent tasks rather than the operative structure of work.
This appendix re-examines all three predictions on samples restricted to the tasks that workers report performing frequently.
The chain-length and neighbor results keep their sign under pruning, though they are estimated less precisely as the samples shrink, and the fragmentation null of Subsection~\ref{sec:fragmentation_prediction} is likewise unchanged.

Rather than re-querying the language model to re-sequence each occupation's surviving tasks, we prune the workflows already generated for the main analysis.
Within each occupation we drop the tasks that fail the frequency filter (discussed below), keep the surviving tasks in their original order, and recompute the average AI chain length, the empirical fragmentation index, and the neighbor-execution flags on the pruned sequence before re-estimating the corresponding statistics. 
Pruning the existing sequences rather than re-prompting for a fresh ordering isolates the question of interest, namely whether infrequent tasks contaminate our measures, while holding the LLM's sequencing logic fixed.

To keep the sample definition consistent across the three predictions, we restrict each cut to the occupations whose frequency-pruned workflow retains at least five tasks, and we evaluate all three predictions on this common set of occupations.
The five-task floor is applied to the pruned workflow, that is, to the steps that survive the frequency filter.
For Predictions~\#1 and~\#3 the chain length and fragmentation index are computed over each occupation's full pruned workflow; for Prediction~\#2 the neighbor regression uses the eligible tasks of the same occupations, namely those with two neighbors on either side, which a five-task workflow always admits, and among these only the tasks whose DWA appears in more than one occupation, as in Section~\ref{sec:DWA_prediction}, so fewer occupations contribute to the neighbor estimates than to the other two.\footnote{For comparability across cuts, the fragmentation regressions z-score all three variables (the AI-execution share, the empirical fragmentation index, and AI exposure) within each cut, so the reported coefficients are standardized betas, where a one-standard-deviation change in fragmentation maps to a $\beta$-standard-deviation change in the AI-execution share. For the neighbor logits we report average marginal effects with analytic robust standard errors clustered at the DWA level, rather than the bootstrap of Table~\ref{tab:DWA_regression_aiExecution_mainSample}. Point estimates for the ``all tasks'' specification of Predictions~\#1 and~\#2 are identical to those reported in the main text. For Prediction~\#3, the five-task floor drops one of the $872$ occupations, which moves the SOC minor-group estimate from $-0.04$ to $-0.05$.}

\subsection{Constructing the Frequent-Task Samples}
\label{app:frequency_robustness_construction}

We measure how often a task is performed using the O*NET \emph{Frequency of Task} (FT) module.
O*NET asks the incumbent workers it surveys in each occupation how often they perform each of the occupation's tasks, offering seven answer categories ordered from least to most frequent: \textit{``yearly or less,'' ``more than yearly,'' ``more than monthly,'' ``more than weekly,'' ``daily,'' ``several times daily,''} and \textit{``hourly or more,''} and it reports the percentage of respondents choosing each category.
A task thus comes not with a single frequency but with a distribution of workers' responses over the seven categories, so restricting the sample to frequently-performed tasks requires two choices: (1) which categories count as ``frequent,'' and (2) how large the share of workers reporting the task in those categories must be.
We keep a task when the combined share of respondents in the designated frequent categories meets a threshold.

We consider three increasingly demanding notions of ``frequent.''
\emph{Daily$+$} counts tasks performed daily, several times daily, or hourly or more; \emph{SeveralDaily$+$} counts those performed several times daily or hourly or more; and \emph{Hourly$+$} counts only those performed hourly or more.
Each notion is evaluated at four thresholds: 20\%, 35\%, 50\%, and 65\% of incumbent worker responses.
A task is retained under the ``Hourly$+$ $\geq 20\%$'' cut, for example, if at least one in five workers reports performing it hourly or more.
To see how the cuts relate to one another, consider a task that 60\% of surveyed workers report performing several times daily, with the remaining responses spread over less frequent categories.
The task survives the SeveralDaily$+$ cuts at the 20\%, 35\%, and 50\% thresholds, and because responses in a stricter category also count toward every looser logic, it survives the corresponding Daily$+$ cuts as well.
Nothing stricter retains it, in either direction. No worker reports the task hourly or more, so it fails every Hourly$+$ cut, and its frequent share falls short of 65\%, so it fails the strictest threshold under both remaining logics.
Each logic-threshold pair therefore defines its own subsample of tasks and occupations.
We estimate every prediction separately on each subsample and present the results for all twelve cuts together in a grid, three logics by four thresholds, for ease of comparison.

Figure~\ref{fig:frequency_pruning_example} makes the pruning procedure concrete for a single occupation, Allergists and Immunologists (SOC 29-1229.01).
\begin{figure}[!t]
\caption{Frequency Pruning of a Single Occupation's Workflow: Allergists and Immunologists}
\label{fig:frequency_pruning_example}
\centering
\begin{subfigure}[b]{\textwidth}
    \centering
    \subcaption{Full Workflow (All Tasks)}
    \includegraphics[width=0.86\textwidth]{plots/taskSequence_vs_anthropicIndex/example_pruning/example_allergists_full.png}
\end{subfigure}

\vspace{0.6em}

\begin{subfigure}[b]{\textwidth}
    \centering
    \subcaption{Pruned Workflow (Hourly$+$ $\geq 50\%$)}
    \includegraphics[width=0.86\textwidth]{plots/taskSequence_vs_anthropicIndex/example_pruning/example_allergists_hourly50.png}
\end{subfigure}
\par\vspace{\fignotesskip}
\begin{minipage}{1\linewidth}
\begin{spacing}{0.2}
{\fontsize{9.5pt}{10pt}\selectfont Notes: Panel~(a) reports the full sequenced workflow for Allergists and Immunologists (SOC 29-1229.01), the object used in the main analysis: the occupation's sixteen O*NET tasks in the order assigned by the language model.
Panel~(b) reports the workflow that survives the Hourly$+$ $\geq 50\%$ frequency cut, which retains a task only if at least half of the surveyed incumbent workers report performing it hourly or more.
Surviving tasks keep their original relative order and their position in the full workflow, so the numbers skipped in Panel~(b) are the pruned steps.
Shaded rows are AI-executed steps (augmentation or automation in the Anthropic Economic Index); unshaded rows are manual.
Pruning keeps all four AI-executed steps but drops the manual steps separating positions~7 and~11, merging two isolated AI-executed steps into a chain of length two.}
\end{spacing}
\end{minipage}
\end{figure}
Panel~(a) shows the occupation's full sequenced workflow, the object used throughout the main analysis: sixteen O*NET tasks in the order the language model places them, with the AI-executed steps shaded.
Panel~(b) shows what survives the Hourly$+$ $\geq 50\%$ cut, one of the most demanding in the grid: seven of the sixteen tasks, retained in their original relative order and still carrying their original position numbers so that the dropped steps can be read off directly.
What the filter removes is what one would expect it to remove.
The specialized diagnostic procedures (steps~3 and~4), the coordination and consultation activities (steps~9 and~10), and the research and continuing-education tasks (steps~14 through~16) are performed episodically rather than within the hour, and they drop out; what remains is the recurring clinical loop of documenting histories, examining patients, interpreting results, planning treatment, educating patients, and prescribing.
The figure also shows why pruning is not innocuous for our measures, which is precisely the concern this appendix addresses.
All four AI-executed steps survive the filter, but in the full workflow each is isolated between manual steps, so the occupation contributes four chains of length one and an average AI chain length of $1.00$; once the intervening manual steps~8 through~10 are dropped, steps~7 and~11 become adjacent and the same four AI-executed tasks now form three chains of lengths $1$, $1$, and $2$, raising the occupation's average to $1.33$.
Pruning can thus both shorten workflows and concentrate the surviving AI-executed steps, and the exercises that follow ask whether our three results survive once this is done for every occupation.
We take the three predictions in turn.


\subsection{Robustness of Prediction \#1 Results}
\label{app:frequency_robustness_chains}

Prediction~\#1 concerns the formation of AI chains.
We measure it, as in Section~\ref{sec:chainLength_prediction}, by the average length of a maximal run of contiguous AI-executed steps, where a step is AI-executed if it is labeled augmented or automated, pooled across all chains in the sample.
A longer-than-random average chain could, however, be delivered mechanically by the composition of each pruned sample.
We therefore benchmark every cut against a within-occupation task position-reshuffle placebo, where we redraw $1{,}000$ reshuffles of task positions within occupations, re-apply the frequency filter, and recompute the statistic, producing a null distribution that holds the composition of the sample fixed while destroying the workflow ordering.
Figure~\ref{fig:chainlength_frequency} reports the observed average against this null as a forest plot: each row is a frequency cut, with the ``all tasks'' baseline analysis in the main manuscript on top; the grey bar spans the 10th to 90th percentiles of the reshuffle null, the vertical tick marks the null mean, and the dot is the observed estimate, red when it falls outside the null band and blue when inside, with the observed percentile printed in the right margin.
\begin{figure}[!t]
\caption{Average AI Chain Length Versus Its Task Position-Reshuffle Null Across Frequency Cuts}
\label{fig:chainlength_frequency}
\centering
\includegraphics[width=0.82\textwidth]{plots/fragmentationIndex_weeklyTasks/chainLength_placebo_forest.png}
\par\vspace{\fignotesskip}
\begin{minipage}{1\linewidth}
\begin{spacing}{0.2}
{\fontsize{9.5pt}{10pt}\selectfont Notes: Each row is a frequency cut, with the ``All tasks'' baseline from the main manuscript on top.
The left labels report the number of occupations in each cut.
The dot is the observed average AI chain length for that cut: the mean length of a maximal run of contiguous AI-executed (augmented or automated) steps.
The grey bar spans the 10th to 90th percentiles of the within-occupation task position-reshuffle null ($1{,}000$ draws), and the vertical tick marks the null mean.
The dot is red if it falls outside the null band and blue if inside.
The right margin reports the observed value and its percentile in the null.}
\end{spacing}
\end{minipage}
\end{figure}
The top all-tasks row reproduces the full-sample value of $1.45$.
Chains shorten modestly as the sample is restricted to frequent tasks, to between $1.24$ and $1.35$ in every cut except the two sparsest Hourly$+$ corners, where they tick back up ($1.40$ and $1.50$, with $75$ and $20$ occupations).
In the full sample the observed value sits at the 100th percentile of its null. AI-executed tasks form longer contiguous chains than random reordering would produce, so the chaining pattern is not a compositional artifact.
Pruning preserves the direction of the result but not its precision.
The observed chain length lies above its null mean in eleven of the twelve pruned cuts and clears the 10-to-90 null band in four of them, all under the inclusive logics, namely Daily$+$ at the $20\%$, $35\%$, and $65\%$ thresholds (the 99th, 100th, and 96th percentiles) and SeveralDaily$+$ at $20\%$ (the 100th).
In the remaining cuts it sits inside the band, at the 58th to 87th percentiles under Daily$+$ and SeveralDaily$+$ and at the 40th to 68th under Hourly$+$, where the samples are smallest and the null bands widest.



\subsection{Robustness of Prediction \#2 Results}
\label{app:frequency_robustness_neighbor}

Figure~\ref{fig:neighbor_frequency_heatmap} repeats the exercise for the neighbor prediction.
\begin{figure}[!t]
\caption{Neighbor AI Execution and Focal-Task AI Execution Across Frequency Cuts}
\label{fig:neighbor_frequency_heatmap}
\centering
\begin{subfigure}[b]{\textwidth}
    \centering
    \subcaption{Previous Task ($k{-}1$)}
    \includegraphics[width=\textwidth]{plots/execTypeVaryingDWA_weeklyTasks/neighbor_logic_threshold_heatmap_prev_bySpec.png}
\end{subfigure}

\vspace{0.8em}

\begin{subfigure}[b]{\textwidth}
    \centering
    \subcaption{Next Task ($k{+}1$)}
    \includegraphics[width=\textwidth]{plots/execTypeVaryingDWA_weeklyTasks/neighbor_logic_threshold_heatmap_next_bySpec.png}
\end{subfigure}
\par\vspace{\fignotesskip}
\begin{minipage}{1\linewidth}
\begin{spacing}{0.2}
{\fontsize{9.5pt}{10pt}\selectfont Notes: Each cell reports the average marginal effect of having the immediately previous task ($k{-}1$, Panel~(a)) or next task ($k{+}1$, Panel~(b)) be AI-executed on the focal task's probability of AI execution, from re-estimating Equation~\eqref{eq:DWA_regression_ai} on the corresponding frequency cut.
The heatmap layout is as in Figure~\ref{fig:frag_frequency}, except that the number beneath each estimate is the number of task observations and the top ``All tasks'' row is the unpruned specification of Section~\ref{sec:DWA_prediction}.
Standard errors are clustered at the DWA level, and significance stars come from the clustered coefficient test. *** p$<$0.01, ** p$<$0.05, * p$<$0.1.
Red denotes a positive effect (the predicted sign) and blue a negative one.
A ``$-$'' marks a cut with too few observations to estimate.}
\end{spacing}
\end{minipage}
\end{figure}
Each cell reports the average marginal effect, from re-estimating Equation~\eqref{eq:DWA_regression_ai}, of having the immediately preceding ($k{-}1$, Panel~(a)) or following ($k{+}1$, Panel~(b)) task be AI-executed on the focal task's probability of AI execution.
A red cell denotes a positive estimate, the sign found in the full sample.
The top row reproduces the full-sample adjacent-step estimates of $+0.11$ for both the previous and the next task under no fixed effects, which attenuate to between $+0.05$ and $+0.06$ once SOC fixed effects absorb cross-family differences.
The estimate is positive in 65 of the 66 estimable pruned cells.
Under the inclusive Daily$+$ logic it is somewhat smaller than in the full sample, at $+0.07$ to $+0.10$ without fixed effects and $+0.03$ to $+0.05$ with SOC fixed effects, and it remains significant in 23 of the 24 cells.
Under the stricter SeveralDaily$+$ and Hourly$+$ logics the estimates are noisier.
The previous-task estimate is largest in the sparsest cuts (SeveralDaily$+$ $\geq 50\%$ and $\geq 65\%$, and Hourly$+$ $\geq 50\%$, with $41$ to $574$ task observations), at $+0.12$ to $+0.33$ and significant throughout, while the next-task estimate in the same cells ranges from $+0.05$ to $+0.17$ and is mostly insignificant.
Outside those cuts, significance is lost mainly in the intermediate SeveralDaily$+$ and Hourly$+$ cuts, and most often with SOC minor-group fixed effects.
The effects of steps two positions away remain small and, under fixed effects, are statistically indistinguishable from zero across the grid, in line with the main text. The part of a step's local context that predicts its execution mode is the part a chain of average length $1.45$ can actually span.\footnote{ To save space, we do not report the results on steps two position away from the focal step and only focus on the outcomes for immediate neighbors.}

Figure~\ref{fig:neighbor_placebo_forest} benchmarks the previous-task ($k{-}1$, Panel~(a)) and next-task ($k{+}1$, Panel~(b)) effects against their position-reshuffle nulls, one column per specification.
\begin{figure}[!t]
\caption{Immediate-Neighbor Associations: Observed Average Marginal Effects Versus Their Position-Reshuffle Nulls Across Frequency Cuts}
\label{fig:neighbor_placebo_forest}
\centering
\begin{subfigure}[b]{\textwidth}
    \centering
    \subcaption{Previous Task ($k{-}1$)}
    \includegraphics[width=\textwidth]{plots/placebo_summary/placebo_summary_forest_neighbor_t1_byCut.png}
\end{subfigure}

\vspace{0.8em}

\begin{subfigure}[b]{\textwidth}
    \centering
    \subcaption{Next Task ($k{+}1$)}
    \includegraphics[width=\textwidth]{plots/placebo_summary/placebo_summary_forest_neighbor_t2_byCut.png}
\end{subfigure}
\par\vspace{\fignotesskip}
\begin{minipage}{1\linewidth}
\begin{spacing}{0.2}
{\fontsize{9.5pt}{10pt}\selectfont Notes: Panel~(a) reports the average marginal effect of having the previous task ($k{-}1$) be AI-executed, and Panel~(b) the corresponding effect for the next task ($k{+}1$).
Rows, the null band, the dot coloring, and the margin labels are as in Figure~\ref{fig:chainlength_frequency}, with the dot now the observed average marginal effect; the three columns within each panel are the no-fixed-effects, SOC major-group, and SOC minor-group specifications.
A blank row denotes a cut with too few neighbored task observations to estimate.}
\end{spacing}
\end{minipage}
\end{figure}
Both estimates sit far in the upper tail of their nulls in the full sample, at the 100th percentile in all three specifications, and mostly stay there under the Daily$+$ logic.
The previous-task estimate lies outside its 10-to-90 null band in ten of the twelve Daily$+$ cells and the next-task estimate in nine, at the 87th to 97th and the 77th to 100th percentiles across those cells, respectively.
Under the stricter logics the two separate.
The previous-task estimate still escapes its null in seven of the twelve SeveralDaily$+$ cells, including all three at the $\geq 50\%$ threshold, but in only one of the nine Hourly$+$ cells, while the next-task estimate stays inside its null band in every SeveralDaily$+$ and Hourly$+$ cell, at the 28th to 86th percentiles.
The positive association between adjacent-step AI execution and the focal step's AI execution thus persists under the inclusive cuts and loses precision, but rarely its sign, under the strictest thresholds, where the surviving task observations number in the low hundreds and the reshuffle bands are correspondingly wide.

\subsection{Robustness of Prediction \#3 Results}
\label{app:frequency_robustness_frag}

Figure~\ref{fig:frag_frequency} summarizes the fragmentation regressions as a heatmap over the logic-by-threshold grid, with the unpruned ``all tasks'' specification from the main analysis in the top row, so that the effect of progressively dropping infrequent tasks can be read down the rows, which impose a more demanding frequency logic, and across the columns, which raise the threshold.
\begin{figure}[!t]
\caption{Empirical Fragmentation Index and AI Execution Across Frequency Cuts}
\label{fig:frag_frequency}
\centering
\includegraphics[width=\textwidth]{plots/efi_matched_exposure/frag_logic_threshold_heatmap_def1_matched.png}
\par\vspace{\fignotesskip}
\begin{minipage}{1\linewidth}
\begin{spacing}{0.2}
{\fontsize{9.5pt}{10pt}\selectfont Notes: Each cell reports the standardized coefficient on the empirical fragmentation index from re-estimating Equation~\eqref{eq:fragmentation_index_regression} on the corresponding frequency cut.
The index is the preferred exposure-based measure.
The number of occupations appears beneath each estimate. No cell in the grid is significant at the 10\% level, so no significance stars are shown. Following the neighbor heatmaps of Subsection~\ref{app:frequency_robustness_neighbor}, the Hourly$+$ and $\geq 65\%$ cut is left blank, since the $20$ occupations it retains are too few for this specification to estimate.
The three heatmaps report the no-fixed-effects, SOC major-group, and SOC minor-group specifications.
Within each heatmap, rows correspond to the three frequency logics and columns to the four thresholds.
The top ``All tasks'' row is the unpruned specification of Section~\ref{sec:fragmentation_prediction}.
Blue denotes a negative coefficient (the predicted sign) and red a positive one. Each panel is scaled to its own largest absolute coefficient, so shading is comparable within a panel but not across panels.}
\end{spacing}
\end{minipage}
\end{figure}
Each cell reports the standardized coefficient on the empirical fragmentation index from re-estimating Equation~\eqref{eq:fragmentation_index_regression} for the corresponding cut, with the number of occupations printed beneath (no cell is significant at the $10\%$ level, so no stars appear); as in the main text, every specification controls for the number of AI-exposed steps in the occupation.
We use our preferred exposure-based measure throughout (as defined in Section~\ref{sec:fragmentation_prediction}, under which both E1- and E2-exposed tasks may form AI chains), and the three heatmaps correspond to the no-fixed-effects, SOC major-group, and SOC minor-group specifications.
A blue cell denotes a negative coefficient, the sign the model predicts, and a red cell a positive one; darker shading indicates a larger magnitude.
The top all-tasks row gives the full-sample standardized estimates of $-0.01$, $-0.09$, and $-0.05$, none of them distinguishable from zero, here estimated on the $871$ occupations that retain at least five tasks and closely matching the headline estimates of Table~\ref{tab:fragmentation_index_regression_exposure}.
Pruning does not recover the relationship.
Across the $33$ pruned cells the coefficient is negative in $13$, the median absolute estimate is $0.05$, and no cell in the grid reaches significance at the $10\%$ level under any of the three specifications.
We do not read the pruned cells as ruling the channel out either, since $20$ of the $33$ retain confidence intervals wide enough to admit a moderate negative effect.
What the grid establishes is that the null of Subsection~\ref{sec:fragmentation_prediction} is a feature of the O*NET sample rather than an artifact of counting rarely-performed tasks as steps.



In short, the chain-length and neighbor patterns keep their sign among frequently performed tasks, and neither they nor the fragmentation null of the O*NET sample appear to be artifacts of the rarely performed activities catalogued in O*NET.